% Chapter 2 -- InitechDOS Commands.  The notation and the information common
% to all commands, then one AN-01 entry per command (IBM 3.10 Reference Ch. 7
% order: alphabetical, one entry per page; docs/pubs/originals/ibm-dos.md
% Sec 8.1).  The entry list is the shipped command table: the internal
% commands of os/milton/command.c cmd_classify (24 names, 18 commands) and
% the external utilities on the shipped disk (Makefile flair-data-rules:
% APPS\CHKDSK.COM ATTRIB.COM LABEL.COM XCOPY.COM TREE.COM).  PATH is a DRAFT
% entry (initech-sy5f: not yet a command).
% Sources (Law 1): os/milton/command.c (cmd_parse, run_external,
% cmd_path_candidates: .COM .EXE .BAT, current directory then PATH),
% os/milton/devices.c (CON AUX PRN CLOCK$ NUL), os/milton/int21.c
% (8.3 search template, '*' and '?').
% GATES: test-command test-shell test-devices test-pubs-screens
\Gates{test-command, test-shell, test-devices, test-pubs-screens}
\chapter{InitechDOS Commands}

\section{Format Notation}
\index{notation}
The Format line of each entry shows exactly how the command is typed.  These
conventions are used:
\begin{itemize}
\item Words in capital letters, such as \kw{COPY}, are keywords.  Type them
  as shown; you may use capital or small letters.
\item Words in \ph{italics} are parameters: replace them with your own
  information, such as a file name.
\item Items in square brackets \opt{\ph{like this}} are optional.  Do not
  type the brackets.
\item A vertical bar between items, as in \kw{ON}\texttt{|}\kw{OFF}, means
  that you type one of them.
\item Three periods \rep{} mean that the item before them may be repeated.
\item Type all other punctuation, such as colons, backslashes, equal signs
  and slashes, exactly as shown.
\end{itemize}

\section{Information Common to All Commands}

\subsection{Entering Commands}
Type a command at the command prompt and press \key{Enter}.  Separate the
command name from its first parameter, and the parameters from each other,
with a space.  A command line may be up to 127 characters long.  InitechDOS
reads the whole line only when you press \key{Enter}; until then you can
correct it with \key{Backspace}, or press \key{Ctrl-C} to cancel it.

\subsection{File Specifications}
\index{file specification}
Commands that work on files take a \emph{file specification}:
\SyntaxLine{\opt{\ph{d}\kw{:}}\opt{\ph{path}}\ph{filename}\opt{\kw{.}\ph{ext}}}
\Param{d:} is the drive.  In this version it is always drive~A, and you may
leave it out.

\Param{path} is the list of directories that leads to the file, each name
followed by a backslash.  A path that begins with a backslash starts at the
root directory; any other path starts at the current directory.  The name
\file{..} stands for the parent of a directory.  If you leave the path out,
InitechDOS uses the current directory.

\Param{filename} is the name of the file, one to eight characters.

\Param{ext} is the extension, one to three characters, separated from the
name by a period.

Make your names from letters and digits.  Do not use spaces, periods
(except the one before the extension), or the characters
\texttt{\textbackslash{} / : * ? " < > | + = ; , [ ]}, which have other
meanings on the command line.  InitechDOS stores and displays names in
capital letters, whatever you type.

\subsection{Wildcard Characters}
\index{wildcards}
Some commands accept the wildcard characters \type{?} and \type{*} in a
file name or extension.  A question mark stands for any one character in
its position.  An asterisk stands for any characters from its position to
the end of the name or of the extension.  \file{*.TXT} names every file
with the extension \file{TXT}; \file{*.*} names every file.  Each entry
states whether the command accepts wildcards.

\subsection{Reserved Device Names}
\index{device names}
These names belong to the devices of the system.  You can use them in place
of a file name, but no file may have one of these names, with or without an
extension:
\par\medskip
\begin{tabular}{@{}l@{\quad}p{80mm}@{}}
{\imSideHeadFont Name} & {\imSideHeadFont Device}\\[2pt]
\cmd{CON}     & The console: the keyboard for input, the display for output\\
\cmd{AUX}     & The first serial port\\
\cmd{PRN}     & The printer\\
\cmd{CLOCK\$} & The system clock\\
\cmd{NUL}     & A dummy device: output to it is discarded, and input from
  it is always empty\\
\end{tabular}
\par\medskip
\index{CON@\texttt{CON}}\index{AUX@\texttt{AUX}}\index{PRN@\texttt{PRN}}%
\index{CLOCK\$@\texttt{CLOCK\$}}\index{NUL@\texttt{NUL}}

\subsection{Internal and External Commands}
\index{internal commands}\index{external commands}
\emph{Internal} commands are part of the command processor,
\file{COMMAND.COM}, and are available at any time:
\par\medskip
\begin{tabular}{@{}llllll@{}}
\cmd{BREAK} & \cmd{CLS}  & \cmd{DEL}   & \cmd{ERASE}  & \cmd{PROMPT} & \cmd{RMDIR}\\
\cmd{CD}    & \cmd{COPY} & \cmd{DIR}   & \cmd{EXIT}   & \cmd{RD}     & \cmd{SET}\\
\cmd{CHDIR} & \cmd{DATE} & \cmd{ECHO}  & \cmd{MD}     & \cmd{REN}    & \cmd{TIME}\\
            &            &             & \cmd{MKDIR}  & \cmd{RENAME} & \cmd{TYPE}\\
            &            &             &              &              & \cmd{VER}\\
            &            &             &              &              & \cmd{VOL}\\
\end{tabular}
\par\medskip
\emph{External} commands are program files on disk.  The external commands
of InitechDOS are \cmd{ATTRIB}, \cmd{CHKDSK}, \cmd{LABEL}, \cmd{TREE} and
\cmd{XCOPY}, kept in the directory \file{\bs APPS} of the InitechDOS
diskette.  Any other program or batch file you run is an external command
too.

When you type a name that is not an internal command, InitechDOS looks for
a file of that name with the extension \file{.COM}, then \file{.EXE}, then
\file{.BAT}, first in the current directory and then in each directory named
in the \cmd{PATH} environment variable, in order.  It runs the first file it
finds.  If the name includes a path, such as \type{\bs APPS\bs CHKDSK},
InitechDOS looks only for that file, with the extension \file{.COM}.  If no
file is found, it displays:
\begin{TypeIn}
Bad command or file name
\end{TypeIn}
The Format line of an external command shows \opt{\ph{d}\kw{:}}%
\opt{\ph{path}} before the command name: the drive and path of the program
file, needed when its directory is not named in \cmd{PATH}.

\subsection{Exit Codes}
\index{exit codes}
When a program ends it returns a number, its \emph{exit code}, to
InitechDOS.  A batch file can test the exit code of the last program it ran
with \cmd{IF ERRORLEVEL} (Chapter~4).  The entries of the external commands
list the codes they return.

\subsection{Messages}
When InitechDOS cannot carry out a command, it displays a message and the
command prompt; nothing is changed.  The messages are listed, with their
meaning and the action to take, in Appendix~A.

% Chapter 2 entries ATTRIB .. DISKCOPY (AN-01).  Sources per entry in the comment
% above it; every Example is a capture listed in screens/SCREENS.

% ATTRIB -- os/milton/dosutil.c attrib_main/attrib_walk (+R -R +A -A /S,
% display R col 1 / A col 4 / full path col 10; messages); shipped as
% APPS\ATTRIB.COM.  Gates: test-attrib, test-dosutil-emu.
\begin{DosEntry}{ATTRIB}{Attribute}{external}
\Purpose Sets or displays the read-only and archive attributes of a file.
\Format \SyntaxLine{\opt{\ph{d}\kw{:}}\opt{\ph{path}}\kw{ATTRIB}
  \opt{\kw{+R}\texttt{|}\kw{-R}} \opt{\kw{+A}\texttt{|}\kw{-A}}
  \opt{\ph{d}\kw{:}}\opt{\ph{path}}\ph{filename}\opt{\kw{.}\ph{ext}}
  \opt{\kw{/S}}}
\Type{External}
\Remarks
\Param{+R} makes the file read-only.  A read-only file can be read, typed
and copied, but it cannot be changed or deleted.

\Param{-R} makes the file read-write again.

\Param{+A} sets the archive attribute of the file.  The archive attribute
marks a file that has been created or changed since it was last backed up;
\cmd{XCOPY} uses it with its \kw{/A} and \kw{/M} parameters.

\Param{-A} clears the archive attribute.

\Param{filename} names the files to change or display.  Wildcards are
allowed.  Directories are not changed.

\Param{/S} processes the matching files in the named directory and in all
of its subdirectories.

If you give no \kw{+} or \kw{-} parameter, \cmd{ATTRIB} displays the
attributes of each matching file: \texttt{R} in the first column for a
read-only file, \texttt{A} in the fourth column when the archive attribute is
set, then the full path of the file.

If no file matches, or the name is a directory, \cmd{ATTRIB} displays
\texttt{Invalid path (or file not found)}.  If you give no file name, or
more than one, it displays \texttt{Invalid number of parameters}.

\cmd{ATTRIB} returns exit code 0 when it succeeds and 1 when it displays a
message.
\Example
The first command makes \file{MEMO.TXT} read-only; the second displays the
attributes of every file with the extension \file{TXT}; the \cmd{DEL}
command then cannot delete the read-only file:
\Shot{attrib}
To display the attributes of the program files in every directory of the
disk, begin at the root directory and use \kw{/S}:
\Shot{attribs}
\end{DosEntry}

% BREAK -- os/milton/command.c builtin_break (ON/OFF case-insensitive, other
% arguments leave the state; "BREAK is on|off"); sysinit.c (default ON
% unless CONFIG.SYS BREAK=OFF); int21.c conin_cooked_line (CTRL-C on CON).
% ADR-0003 Amendment DEC-16.  Gates: test-command, test-j1lu-control-emu.
\begin{DosEntry}{BREAK}{Control-Break}{internal}
\Purpose Sets or displays the CTRL-BREAK check.
\Format \SyntaxLine{\kw{BREAK} \opt{\kw{ON}\texttt{|}\kw{OFF}}}
\Type{Internal}
\Remarks
Type \cmd{BREAK} with no parameter to display the current setting:
\texttt{BREAK is on} or \texttt{BREAK is off}.  After \kw{ON} or \kw{OFF},
\cmd{BREAK} also displays the new setting.

The check is on when InitechDOS starts, unless \file{CONFIG.SYS} contains
the configuration command \texttt{BREAK=OFF} (Chapter~3).  Any parameter
other than \kw{ON} or \kw{OFF} leaves the setting unchanged.

Whatever the setting, InitechDOS checks for \key{Ctrl-C} whenever it reads
a line from the keyboard: the line is discarded, a program reading it is
stopped, and the command prompt is displayed again.
\draft{see bead initech-vlhp (DEC-16 C-6): in DOS 3.3, BREAK ON also checks for CTRL-C during disk
and other operations; the shipped build records and reports the setting
but checks only during keyboard input}
With \kw{ON}, InitechDOS also checks for \key{Ctrl-C} while a program reads
or writes the disk, so a long operation can be stopped.
\Example
\Shot{break}
\end{DosEntry}

% CD -- os/milton/command.c builtin_cd (no arg: cwd_display; failure:
% MSG-DOS-0018 "Invalid directory"); int21.c do_chdir (relative, absolute,
% "..", "\" alone = the root since initech-33t2).  Gates: test-ut6d,
% test-fat-subdir, test-shell.
\begin{DosEntry}{CD}{Change Directory}{internal}
\AlsoNamed{CHDIR}
\Purpose Changes the current directory, or displays its name.
\Format \SyntaxLine{\kw{CD} \opt{\ph{d}\kw{:}}\opt{\ph{path}}}
  \SyntaxLine{\kw{CHDIR} \opt{\ph{d}\kw{:}}\opt{\ph{path}}}
\Type{Internal}
\Remarks
\Param{path} names the directory that is to become the current directory.
A path that begins with a backslash starts at the root directory; any other
path starts at the current directory.  \file{..} names the parent of the
current directory.

Type \cmd{CD} with no parameter to display the drive and full path of the
current directory.

The command prompt shows the current directory, so a successful \cmd{CD}
displays nothing more.  If the directory does not exist, \cmd{CD} displays
\texttt{Invalid directory} and the current directory does not change.

\type{CD \bs} makes the root directory the current directory, from any
directory on the disk.
\Example
The example moves into \file{SAMPLES}, displays the current directory,
moves back with \type{CD ..}, moves in again and goes straight to the root
with \type{CD \bs}:
\Shot{cd}
\end{DosEntry}

% CHKDSK -- os/milton/chkdsk.c (parse: one file operand, /F /V; report
% lines; contiguity; lost chains + FILE0000.CHK; cross-links reported, never
% repaired; memory from the live MCB chain); chkdsk_program.c.  Shipped as
% DOS\CHKDSK.COM.  Gates: test-chkdsk, test-chkdsk-emu.
\begin{DosEntry}{CHKDSK}{Check Disk}{external}
\Purpose Checks the directories and the file allocation table of a disk and
reports on the disk and on memory.
\Format \SyntaxLine{\opt{\ph{d}\kw{:}}\opt{\ph{path}}\kw{CHKDSK}
  \opt{\ph{d}\kw{:}}\opt{\ph{path}}\opt{\ph{filename}\opt{\kw{.}\ph{ext}}}
  \opt{\kw{/F}} \opt{\kw{/V}}}
\Type{External}
\Remarks
\cmd{CHKDSK} reads every directory of the disk and follows the chain of
allocation units of every file.  It then displays a status report:
\begin{itemize}
\item the volume label and the date and time the volume was created
\item the total disk space, in bytes
\item the bytes and number of hidden files, of directories, and of user
  files
\item the bytes still available on the disk
\item the total memory and the free memory, in bytes
\end{itemize}

\Param{filename} names files to check for fragmentation.  Wildcards are
allowed.  For each named file stored in more than one piece, \cmd{CHKDSK}
displays its path and \texttt{Contains} \ph{n} \texttt{non-contiguous
blocks}; when every named file is in one piece it displays \texttt{All
specified file(s) are contiguous}.  If no file matches, it displays
\texttt{File not found}.

\Param{/F} lets \cmd{CHKDSK} correct the errors it finds.  Without
\kw{/F}, \cmd{CHKDSK} reports what it would do and displays
\texttt{Corrections will not be written to disk}; nothing on the disk is
changed.

\Param{/V} displays the full path of every file as it is checked.

When \cmd{CHKDSK} finds allocation units that are marked in use but belong
to no file, it displays
\ph{n} \texttt{lost allocation units found in} \ph{n} \texttt{chains.}
and asks \texttt{Convert lost chains to files (Y/N)?}  With \kw{/F}, answer
\type{Y} to save each chain as a file named \file{FILE0000.CHK},
\file{FILE0001.CHK}, and so on, in the root directory, where you can examine
it; answer \type{N} to free the space.  Without \kw{/F}, \cmd{CHKDSK}
displays the number of bytes that would be freed.

A file whose allocation chain points outside the disk is reported with
\texttt{Has invalid cluster, file truncated}.  Two files that share an
allocation unit are both reported with \texttt{File is cross-linked: on
cluster} \ph{n}; \cmd{CHKDSK} does not repair a cross-link.  Copy both
files, check their contents, and delete the damaged one.

\cmd{CHKDSK} returns exit code 0 when the disk has no errors and 1 when it
reports errors.
\Example
To check the practice diskette:
\Shot{chkdsk}
To check whether \file{MEMO.TXT} is stored in one piece:
\Shot{chkfile}
\end{DosEntry}

% CLS -- os/milton/command.c builtin_cls (form feed: clear and home; any
% parameters ignored).  Gate: test-command, test-pubs-screens (every screen
% begins with CLS).
\begin{DosEntry}{CLS}{Clear Screen}{internal}
\Purpose Clears the display.
\Format \SyntaxLine{\kw{CLS}}
\Type{Internal}
\Remarks
The display is cleared and the command prompt appears on the top line.
Parameters typed after \cmd{CLS} are ignored.
\Example
Type \type{CLS} and press \key{Enter}.  The screen is cleared and shows only
the command prompt, \texttt{A:\bs>}.  Every example in this book begins
this way.
\end{DosEntry}

% COPY -- os/milton/command.c builtin_copy (two operands; directory target
% keeps the name; same-file refusal MSG-DOS-0020 + "0 File(s) copied";
% source missing MSG-DOS-0003; create refused MSG-DOS-0009 or MSG-DOS-0002;
% footer "        1 File(s) copied", MSG-DOS-0061 since WL-0154); COPY CON
% ends at Ctrl-Z (WL-0102).
% Gates: test-copy-selfcopy, test-dos-safety-k01, test-dos-safety-k02,
% test-j1lu-control-emu.  Date stamp of the copy: initech-dwbl.
\begin{DosEntry}{COPY}{Copy File}{internal}
\Purpose Copies a file.
\Format \SyntaxLine{\kw{COPY} \opt{\ph{d}\kw{:}}\opt{\ph{path}}\ph{filename}\opt{\kw{.}\ph{ext}}
  \opt{\ph{d}\kw{:}}\opt{\ph{path}}\opt{\ph{filename}\opt{\kw{.}\ph{ext}}}}
\Type{Internal}
\Remarks
The first file specification names the \emph{source}, the file to copy.
The second names the \emph{target}: the new file, or the directory to copy
into.  Both are required.

If the target is a directory, the copy is put in that directory with the
name of the source file.  If the target is an existing file, it is
replaced.  When the copy is complete, \cmd{COPY} displays
\texttt{1 File(s) copied}.\draft{see bead initech-bc2j: the shipped COPY
also accepts wildcards, combines files named with +, takes the /A, /B and
/V switches, and copies into the current directory when no target is given;
this entry is to be rewritten}

To copy a directory and its subdirectories, use \cmd{XCOPY}.

You can copy from the keyboard by naming the device \file{CON} as the
source: type the lines of the file, then press \key{Ctrl-Z} and
\key{Enter}.  \texttt{COPY CON NOTE.TXT} creates the file \file{NOTE.TXT}
from what you type.

If the source and target are the same file, \cmd{COPY} displays
\texttt{File cannot be copied onto itself} and \texttt{0 File(s) copied},
and the file is not changed.  If the source does not exist, it displays
\texttt{File not found}.  If the target is read-only, it displays
\texttt{Access denied}; if the target cannot be created for another reason,
such as a full disk, it displays \texttt{Bad command or file name}.

\draft{bead needed (follow-up of initech-dwbl): the shipped COPY stamps the
copy with the current date and time; DOS 3.3 keeps the source's date and
time}
The copy has the same date and time as the source file.
\Example
The first command copies \file{MEMO.TXT} to a new file, \file{MEMO.BAK}; the
second copies \file{PHONE.TXT} into the directory \file{SAMPLES}; the third
names the same file twice:
\Shot{copy}
\end{DosEntry}

% DATE -- os/milton/command.c builtin_date, cmd_parse_date (mm-dd-yy[yy],
% '-' or '/' used twice; 80..99 -> 19xx, 00..79 -> 20xx; four-digit years from 1980 on, upper bound in cmd_parse_date;
% days per month with leap years), cmd_format_date ("Current date is Day
% mm-dd-yyyy"); "Invalid date" then the unchanged date.  No prompt for a new
% date in the shipped build (documented as built).  Gates: test-datetime,
% test-command.
\begin{DosEntry}{DATE}{}{internal}
\Purpose Sets or displays the system date.
\Format \SyntaxLine{\kw{DATE} \opt{\ph{mm}\kw{-}\ph{dd}\kw{-}\ph{yy}}}
\Type{Internal}
\Remarks
\Param{mm} is the month, 1 through 12.

\Param{dd} is the day, 1 through the last day of the month.

\Param{yy} is the year.  Two digits from 80 through 99 stand for the years
from 1980 to the end of the century; two digits from 00 through 79 stand for
the years of the next century.  You may also give all four digits.

Separate the parts with hyphens or with slashes, the same character both
times.  InitechDOS works out the day of the week for you.

With a valid date, \cmd{DATE} sets the system date and displays it.  With
no parameter, it displays the current date:
\texttt{Current date is} followed by the day of the week and the date.

If the date is not valid (for example, February~30), \cmd{DATE} displays
\texttt{Invalid date}, leaves the system date unchanged and displays it.

\Example
\Shot{date}
\end{DosEntry}

% DEL -- os/milton/command.c builtin_del (wildcards: restart FINDFIRST per
% unlink; *.* asks MSG-DOS-0012 first, only Y deletes; read-only:
% MSG-DOS-0009 + the names left; MSG-DOS-0003; MSG-DOS-0011).  Gates:
% test-dos-safety-k03, test-dos-safety-k04, test-dos-safety-k15,
% test-dos-safety-del-failure.
\begin{DosEntry}{DEL}{Delete}{internal}
\AlsoNamed{ERASE}
\Purpose Deletes files.
\Format \SyntaxLine{\kw{DEL} \opt{\ph{d}\kw{:}}\opt{\ph{path}}\ph{filename}\opt{\kw{.}\ph{ext}}}
  \SyntaxLine{\kw{ERASE} \opt{\ph{d}\kw{:}}\opt{\ph{path}}\ph{filename}\opt{\kw{.}\ph{ext}}}
\Type{Internal}
\Remarks
\Param{filename} names the file or files to delete.  Wildcards are
allowed: \type{DEL *.BAK} deletes every file with the extension \file{BAK}
in the current directory.

To delete every file in a directory, name \file{*.*}.  \cmd{DEL} first asks
\texttt{Are you sure (Y/N)?}  Type \type{Y} and press \key{Enter} to delete
the files; any other answer deletes nothing.

\cmd{DEL} displays nothing when it succeeds.  If no file matches, it
displays \texttt{File not found}.  A read-only file is not deleted:
\cmd{DEL} displays \texttt{Access denied}, and when you used wildcards it
also lists the matching files that remain.  \cmd{DEL} does not delete directories; use \cmd{RD}.

\begin{Important}
A deleted file cannot be recovered with InitechDOS commands.
\end{Important}
\Example
The example makes a copy of \file{MEMO.TXT}, refuses to delete every file,
deletes the copy, and shows that it is gone:
\Shot{del}
\end{DosEntry}

% DIR -- os/milton/command.c builtin_dir (MSG-DOS-0021..0026 display
% fragments, cmd_format_dir_line, footer from AH=36h; empty -> File not
% found; e47v: one operand -- drive A only, else MSG-DOS-0004 -- a directory
% or a file name, "name" = name.* and ".ext" = *.ext, /P pauses with
% MSG-DOS-0016 after 23 rows, /W five names to a line, a subdirectory in
% brackets; any other switch or a second operand MSG-DOS-0010).
% Gates: test-dir, test-dir-look, test-dir-volume, test-dir-free-space,
% test-dir-page, test-dir-args-mutant (initech-e47v).
\begin{DosEntry}{DIR}{Directory}{internal}
\Purpose Lists the files and subdirectories in a directory.
\Format \SyntaxLine{\kw{DIR} \opt{\ph{d}\kw{:}}\opt{\ph{path}}\opt{\ph{filename}\opt{\kw{.}\ph{ext}}}
  \opt{\kw{/P}} \opt{\kw{/W}}}
\Type{Internal}
\Remarks
\Param{d:} is the drive; it can only be drive~A.  Any other drive displays
\texttt{Invalid drive specification}.

\Param{path} names the directory to list.  If you leave it out, \cmd{DIR}
lists the current directory.  Listing a directory does not change the
current directory.

\Param{filename} limits the listing to matching files.  Wildcards are
allowed.  A name without an extension lists every extension of that name
(\type{DIR MEMO} lists \file{MEMO.*}), and an extension alone lists every
name with that extension (\type{DIR .TXT} lists \file{*.TXT}).

\Param{/P} pauses when the screen is full and displays
\texttt{Press any key to continue . . .}; press any key to see the next
screen.

\Param{/W} lists the names only, five to a line, with the name of a
subdirectory in brackets.

Only one file specification may be given; a second one, or any switch
other than \kw{/P} and \kw{/W}, displays \texttt{Too many parameters}.

The listing begins with the volume label of the disk and the name of the
directory.  Each file then has one line: its name, its extension, its size
in bytes, and the date and time it was last changed, with \texttt{a} for
morning and \texttt{p} for afternoon.  A subdirectory shows
\texttt{<DIR>} in place of a size.  The last line gives the number of
entries listed and the number of bytes free on the disk.

If the directory has no entries, \cmd{DIR} displays the headings followed
by \texttt{File not found}.
\Example
To list the root directory of the practice diskette:
\Shot{dir}
To list only the text files, and then the directory \file{SAMPLES} in the
wide form:
\Shot{dirargs}
\end{DosEntry}

% DISKCOMP / DISKCOPY -- os/milton/diskutil.c copy_run (the single-drive
% swap flow of MS-DOS 3.3 UG p.26 and pp.59-63: FIRST/SECOND and
% SOURCE/TARGET prompts, 64 sectors per pass; "Compare error on disk / side
% s, track t", "Compare OK"; removable drives only, two-line refusal p.358;
% B: is MSG-DOS-0004 by DEC-18, spec/drive_letters.json; /1 /8 options;
% "Copy/Compare another diskette (Y/N)?").  WL-0147.  Shipped as
% DOS\DISKCOMP.COM and DOS\DISKCOPY.COM.  Gates: test-formatsys,
% test-formatsys-emu, test-diskcopy-emu.
\begin{DosEntry}{DISKCOMP}{Compare Diskette}{external}
\Purpose Compares the contents of two diskettes sector by sector.
\Format \SyntaxLine{\opt{\ph{d}\kw{:}}\opt{\ph{path}}\kw{DISKCOMP}
  \opt{\ph{d}\kw{:} \opt{\ph{d}\kw{:}}} \opt{\kw{/1}} \opt{\kw{/8}}}
\Type{External}
\Remarks
\Param{d:} names the drives of the first and second diskette.  If you give
none, or only one, drive~A is used for the others.  Both drives must be
diskette drives: for a fixed disk \cmd{DISKCOMP} displays
\texttt{Specified drive does not exist,} \texttt{or is non-removable}.  Drive~B
does not exist in this version (\texttt{Invalid drive specification}).

With one drive, \cmd{DISKCOMP} asks for the diskettes in turn:
\texttt{Insert FIRST diskette into drive A:} and \texttt{Insert SECOND
diskette into drive A:}.  Change the diskette and press any key.  It reads
up to 64 sectors of the first diskette at a time, so a large diskette is
compared in several passes.

\Param{/1} compares only the first side of each diskette.  \Param{/8}
compares only the first eight sectors of each track.

For every sector that differs, \cmd{DISKCOMP} displays
\texttt{Compare error on disk} and, on the next line, the side and track.
If no sector differs, it displays \texttt{Compare OK}.  It then asks
\texttt{Compare another diskette (Y/N)?}; type \type{Y} or \type{N}.
\Example
To compare two diskettes in drive~A, type \type{DISKCOMP A: A:}, then
insert the first and second diskettes as \cmd{DISKCOMP} asks for them.
\end{DosEntry}

\begin{DosEntry}{DISKCOPY}{Copy Diskette}{external}
\Purpose Copies the contents of one diskette onto another, sector by sector.
\Format \SyntaxLine{\opt{\ph{d}\kw{:}}\opt{\ph{path}}\kw{DISKCOPY}
  \opt{\ph{d}\kw{:} \opt{\ph{d}\kw{:}}} \opt{\kw{/1}}}
\Type{External}
\Remarks
\Param{d:} names the drives of the source and target diskette.  If you give
none, or only one, drive~A is used for the others.  Both drives must be
diskette drives: for a fixed disk \cmd{DISKCOPY} displays
\texttt{Specified drive does not exist,} \texttt{or is non-removable}.  Drive~B
does not exist in this version (\texttt{Invalid drive specification}).

With one drive, \cmd{DISKCOPY} asks for the diskettes in turn:
\texttt{Insert SOURCE diskette in drive A:} and \texttt{Insert TARGET
diskette into drive A:}, each followed by \texttt{Press any key when ready
. . .}  It reads up to 64 sectors of the source at a time, so a large
diskette is copied in several passes.  Everything on the target diskette is
replaced.

\Param{/1} copies only the first side of the diskette.

If the source cannot be read, \cmd{DISKCOPY} displays \texttt{FIRST
diskette bad or incompatible}; if the target cannot be written, it displays
\texttt{Target diskette may be unusable}.  When the copy is complete it asks
\texttt{Copy another diskette (Y/N)?}; type \type{Y} or \type{N}.
\Example
To copy a diskette using drive~A only, type \type{DISKCOPY A: A:}, then
insert the source and target diskettes as \cmd{DISKCOPY} asks for them.
\end{DosEntry}

% Chapter 2 entries ECHO .. REN (AN-01).

% ECHO -- os/milton/command.c builtin_echo (no arg: "ECHO is on"; else the
% text as typed).  ECHO ON/OFF in batch files: Chapter 4 (batch.c).
% Gates: test-command, test-batch.
\begin{DosEntry}{ECHO}{}{internal}
\Purpose Displays a message, or the state of batch-file echoing.
\Format \SyntaxLine{\kw{ECHO} \opt{\ph{message}}}
\Type{Internal}
\Remarks
\Param{message} is the text to display.  It is displayed exactly as you
type it, in capital and small letters.

Typed with no message, \cmd{ECHO} displays \texttt{ECHO is on}.

In a batch file, \cmd{ECHO ON} and \cmd{ECHO OFF} control whether the
command lines of the batch file are displayed as they are carried out.  See
``ECHO Subcommand'' in Chapter~4.
\Example
\Shot{echo}
\end{DosEntry}

% EXIT -- os/milton/command.c dispatch_line CMD_EXIT (a permanent processor
% ignores EXIT without a message; a secondary returns to its parent).
% Gates: test-shell-permanent, test-exit-handles.
\begin{DosEntry}{EXIT}{}{internal}
\Purpose Ends a secondary command processor and returns to the program that
started it.
\Format \SyntaxLine{\kw{EXIT}}
\Type{Internal}
\Remarks
A program can start a second copy of the command processor so that you can
type commands without leaving the program.  \cmd{EXIT} ends that copy and
returns you to the program.

The command processor that InitechDOS starts when the system is turned on
is permanent.  At its prompt \cmd{EXIT} has no effect and displays nothing.
\Example
At the \texttt{A:\bs>} prompt of a freshly started system, type
\type{EXIT} and press \key{Enter}.  The prompt is displayed again, because
the command processor started with the system is permanent.
\end{DosEntry}

% LABEL -- os/milton/dosutil.c label_main / label_read / label_write
% (current drive only; prompts; the invalid characters; 11 characters,
% capitals; root VOLLABEL entry + extended BPB label field; messages).
% Shipped as APPS\LABEL.COM.  Gates: test-label-vol, test-dosutil-emu.
\begin{DosEntry}{LABEL}{Volume Label}{external}
\Purpose Creates, changes or deletes the volume label of a disk.
\Format \SyntaxLine{\opt{\ph{d}\kw{:}}\opt{\ph{path}}\kw{LABEL}
  \opt{\ph{d}\kw{:}}\opt{\ph{label}}}
\Type{External}
\Remarks
\Param{d:} is the drive of the disk.  Only the current drive can be
labeled; another drive letter displays \texttt{Invalid drive
specification}.

\Param{label} is the new volume label, up to 11 characters.  Letters are
changed to capitals; characters after the eleventh are ignored.  A label
may contain spaces, but not the characters
\texttt{* ? / \textbackslash{} | . , ; : + = < > [ ] ( ) \& \textasciicircum{} "};
\cmd{LABEL} displays \texttt{Invalid characters in volume label} and
leaves the label unchanged.

If you type \cmd{LABEL} with no label, it displays the current label (or
\texttt{has no label}) and asks for a new one:
\begin{TypeIn}
Volume label (11 characters, ENTER for none)?
\end{TypeIn}
Type the new label and press \key{Enter}.  If you press \key{Enter} alone
and the disk has a label, \cmd{LABEL} asks
\texttt{Delete current volume label (Y/N)?}  Type \type{Y} to delete the
label; any other answer keeps it.

The label is shown by \cmd{DIR}, \cmd{VOL}, \cmd{CHKDSK} and \cmd{TREE}.
If the root directory has no room for the label, \cmd{LABEL} displays
\texttt{No room in root directory}.
\Example
To give the practice diskette the label \file{REPORTS}, answering the
question, and display it:
\Shot{label}
To set the label in one command, and an attempt with a period in it:
\Shot{labelset}
\end{DosEntry}

% MD -- os/milton/command.c builtin_md (MSG-DOS-0011, MSG-DOS-0017; silent
% success); int21.c AH=39h.  Directory date stamp: initech-dwbl.
% Gates: test-ut6d, test-fat12-mkdir, test-fat-subdir.
\begin{DosEntry}{MD}{Make Directory}{internal}
\AlsoNamed{MKDIR}
\Purpose Creates a subdirectory.
\Format \SyntaxLine{\kw{MD} \opt{\ph{d}\kw{:}}\ph{path}}
  \SyntaxLine{\kw{MKDIR} \opt{\ph{d}\kw{:}}\ph{path}}
\Type{Internal}
\Remarks
\Param{path} names the new directory.  A path that begins with a backslash
starts at the root directory; any other path starts at the current
directory.  Every directory named before the last one must already exist.
A directory name has the same form as a file name.

\cmd{MD} displays nothing when it succeeds.  A new directory contains two
entries: \file{.}, the directory itself, and \file{..}, its parent.

If the directory cannot be created (for example, because a file or
directory of that name already exists), \cmd{MD} displays
\texttt{Unable to create directory}.  With no path it displays
\texttt{Required parameter missing}.

\draft{initech-dwbl: the shipped build stamps new directories 0-00-80
12:00a (see the example); DOS 3.3 stamps them with the current date and
time}
The new directory is stamped with the current date and time.
\Example
To create the directory \file{REPORTS}, make it the current directory and
list it:
\Shot{md}
\end{DosEntry}

% PATH -- not a command in the shipped build (cmd_classify has no PATH row;
% typing it displays "Bad command or file name").  initech-sy5f.  The PATH
% environment variable itself IS used (command.c run_external,
% cmd_path_candidates; SET PATH= gated by test-command, test-env).
\begin{DosEntry}[initech-sy5f: PATH is not yet an internal command in the
shipped build; the entry describes DOS 3.3.  SET PATH= works today]{PATH}{}{internal}
\Purpose Sets or displays the list of directories searched for commands.
\Format \SyntaxLine{\kw{PATH} \opt{\opt{\ph{d}\kw{:}}\ph{path}\opt{\kw{;}\opt{\ph{d}\kw{:}}\ph{path}}\rep}}
  \SyntaxLine{\kw{PATH} \kw{;}}
\Type{Internal}
\Remarks
\Param{path} names a directory to search.  Separate the directories with
semicolons.  When a command is not an internal command and is not in the
current directory, InitechDOS searches these directories, in order.

Typed with no parameter, \cmd{PATH} displays the current list, as
\texttt{PATH=}\ph{list}.  \type{PATH ;} clears the list, so that only the
current directory is searched.

\cmd{PATH} stores the list in the environment variable \texttt{PATH}; you
can also set it with \type{SET PATH=A:\bs APPS}, as the practice
diskette's \file{AUTOEXEC.BAT} does.
\Example
\type{PATH A:\bs APPS} makes InitechDOS search the directory
\file{\bs APPS} for commands.
\end{DosEntry}

% PROMPT -- os/milton/command.c builtin_prompt (empty -> $P$G; env PROMPT;
% "Out of environment space"), cmd_render_prompt ($P $G $L $B $Q $N $T $D
% $$ $_; any other $x prints x).  Defects found: $D shows 00-00-0000 (the
% REPL passes no date) and $$ cuts the prompt (AH=09h '$' terminator) --
% beads needed, WL-0119.  Gates: test-command, test-shell.
\begin{DosEntry}{PROMPT}{}{internal}
\Purpose Changes the command prompt.
\Format \SyntaxLine{\kw{PROMPT} \opt{\ph{prompt-text}}}
\Type{Internal}
\Remarks
\Param{prompt-text} is the new prompt.  Characters are displayed as you
type them, except that a dollar sign and the character after it stand for
the following:
\par\medskip
\begin{tabular}{@{}l@{\quad}p{75mm}@{}}
{\imSideHeadFont Code} & {\imSideHeadFont Displays}\\[2pt]
\kw{\$P} & The current drive and directory, for example \texttt{A:\bs SAMPLES}\\
\kw{\$G} & The \texttt{>} character\\
\kw{\$L} & The \texttt{<} character\\
\kw{\$B} & The \texttt{|} character\\
\kw{\$Q} & The \texttt{=} character\\
\kw{\$N} & The current drive letter\\
\kw{\$T} & The time, as \texttt{hh:mm:ss.cc}\\
\kw{\$\_} & Starts a new line\\
\end{tabular}
\par\medskip
A dollar sign followed by any other character displays that character.

\draft{bead needed: in the shipped build \$D displays 00-00-0000 and \$\$
ends the prompt at that point; DOS 3.3 displays the date and a dollar sign}
\kw{\$D} displays the date, and \kw{\$\$} displays the \texttt{\$}
character.

Typed with no parameter, \cmd{PROMPT} restores the standard prompt,
\kw{\$P\$G}.  The prompt is kept in the environment variable
\texttt{PROMPT}, which \cmd{SET} displays.  If there is no room for it,
\cmd{PROMPT} displays \texttt{Out of environment space}.
\Example
The first command makes the prompt the drive letter alone; the second makes
it the word \texttt{Ready}; the third restores the standard prompt:
\Shot{prompt}
\end{DosEntry}

% RD -- os/milton/command.c builtin_rd (MSG-DOS-0011, MSG-DOS-0019; silent
% success).  Gates: test-ut6d, test-fat-subdir.
\begin{DosEntry}{RD}{Remove Directory}{internal}
\AlsoNamed{RMDIR}
\Purpose Removes a subdirectory.
\Format \SyntaxLine{\kw{RD} \opt{\ph{d}\kw{:}}\ph{path}}
  \SyntaxLine{\kw{RMDIR} \opt{\ph{d}\kw{:}}\ph{path}}
\Type{Internal}
\Remarks
\Param{path} names the directory to remove.  The directory must be empty
except for its \file{.} and \file{..} entries.  You cannot remove the root
directory or the current directory.

\cmd{RD} displays nothing when it succeeds.  If the directory does not
exist, is not empty, or is the current directory, \cmd{RD} displays
\texttt{Invalid path, not directory, or directory not empty} and nothing is
removed.  Delete the files in a directory with \cmd{DEL} before you remove
it.
\Example
The example creates and removes the directory \file{TEMP}, then tries to
remove \file{SAMPLES}, which still holds files:
\Shot{rd}
\end{DosEntry}

% REN -- os/milton/command.c builtin_ren (AH=56h, same directory; failure
% MSG-DOS-0002; MSG-DOS-0011; silent success).  Gates: test-command,
% test-fat12-dir.
\begin{DosEntry}{REN}{Rename}{internal}
\AlsoNamed{RENAME}
\Purpose Changes the name of a file.
\Format \SyntaxLine{\kw{REN} \opt{\ph{d}\kw{:}}\opt{\ph{path}}\ph{filename}\opt{\kw{.}\ph{ext}}
  \ph{filename}\opt{\kw{.}\ph{ext}}}
  \SyntaxLine{\kw{RENAME} \opt{\ph{d}\kw{:}}\opt{\ph{path}}\ph{filename}\opt{\kw{.}\ph{ext}}
  \ph{filename}\opt{\kw{.}\ph{ext}}}
\Type{Internal}
\Remarks
The first file specification names the file to rename.  The second is the
new name; do not give it a drive or path.  The file stays in its directory.
Wildcards are not allowed in this version.

\cmd{REN} displays nothing when it succeeds.  If the file does not exist,
or a file with the new name already exists, \cmd{REN} displays
\texttt{Bad command or file name}.  With only one name it displays
\texttt{Required parameter missing}.
\Example
To rename \file{MEMO.TXT} to \file{MEMO.OLD} and list the directory:
\Shot{ren}
\end{DosEntry}

% Chapter 2 entries SET .. XCOPY (AN-01).

% SET -- os/milton/command.c builtin_set, cmd_set_parse (list; NAME=VALUE,
% name upcased; NAME= clears, silent; NAME alone queries:
% "Environment variable not defined"); os/milton/env.c (512-byte arena:
% "Out of environment space"); command_repl (COMSPEC, PROMPT, PATH).
% Gates: test-env, test-command, test-interp-set.
\begin{DosEntry}{SET}{}{internal}
\Purpose Sets, removes or displays environment variables.
\Format \SyntaxLine{\kw{SET} \opt{\ph{name}\opt{\kw{=}\opt{\ph{string}}}}}
\Type{Internal}
\Remarks
The \emph{environment} is an area of memory that holds named strings, called
environment variables, for InitechDOS and for the programs you run.

Typed with no parameter, \cmd{SET} lists every variable, one to a line, as
\ph{name}\texttt{=}\ph{string}.

\Param{name=string} sets the variable \ph{name} to \ph{string}, replacing
any earlier value.  The name is changed to capital letters; the string is
kept exactly as you type it.

\Param{name=} with no string removes the variable.  Nothing is displayed.

\Param{name} alone displays that variable, or
\texttt{Environment variable not defined}.

InitechDOS sets three variables when it starts: \texttt{COMSPEC}, the
location of the command processor; \texttt{PROMPT}, the command prompt (see
\cmd{PROMPT}); and \texttt{PATH}, the directories searched for commands
(see Chapter~1, ``The Practice Diskette'').  A batch file can use the value
of a variable by writing its name between percent signs, as
\texttt{\%PATH\%} (Chapter~4).

The environment holds 512 characters.  If a new variable does not fit,
\cmd{SET} displays \texttt{Out of environment space} and the environment is
not changed.
\Example
To list the environment, add the variable \texttt{DEPT}, and list it
again:
\Shot{set}
To display one variable, remove it, and look for it again:
\Shot{setclr}
\end{DosEntry}

% initech-2kgh: sort_program.asm, existing m0dc ascending/reverse goldens.
% Ref: MS-DOS 3.3 User's Reference p.116/PDF201.
\Gates{test-sort-filter, test-sort-filter-mutant, test-dos-folder}
\begin{DosEntry}{SORT}{}{external}
\Purpose Sorts input lines and displays them in order.
\Format \SyntaxLine{\kw{SORT} \opt{\kw{/R}} \kw{<} \ph{filename}}
\Type{External}
\Remarks
Input may come from a file or a pipe. Letters are sorted without regard
to capital or small letters. \Param{/R} reverses the order. Use \type{>}
and a file name to save the sorted output.
\Example
To save sorted lines, type \type{SORT < PHONE.TXT > SORTED.TXT}.
\end{DosEntry}

% SYS -- os/milton/diskutil.c sys_run (MS-DOS 3.3 UG p.119: the two system
% files, in order, into the first two root entries, fragmented sources
% allowed, COMMAND.COM NOT copied; the destination BPB kept, the InitechDOS
% boot code written; any failure "No room for system on destination
% disk"; success "System transferred").  Source = drive C (DEC-18).
% WL-0147.  Shipped as DOS\SYS.COM.  Gates: test-formatsys,
% test-formatsys-emu, test-format-system-boot.
\begin{DosEntry}{SYS}{System}{external}
\Purpose Transfers the InitechDOS system files to a disk.
\Format \SyntaxLine{\opt{\ph{d}\kw{:}}\opt{\ph{path}}\kw{SYS} \ph{d}\kw{:}}
\Type{External}
\Remarks
\Param{d:} is the drive of the disk to receive the system; it is required.
\cmd{SYS} copies the two system files, \file{IO.SYS} and \file{INITDOS.SYS},
from drive~C to the first two entries of the disk's root directory, and
writes the InitechDOS start-up code into its boot record.  The files already
on the disk are kept.  \cmd{SYS} does not copy \file{COMMAND.COM}; copy it
with \cmd{COPY}.

The first two root directory entries must be free, or already hold the
system files, as on a disk formatted with \kw{FORMAT /B}.  When the system
has been transferred, \cmd{SYS} displays \texttt{System transferred}.  If
the entries are in use, the disk has not enough free space, or the system
files cannot be read, it displays \texttt{No room for system on destination
disk} and the disk is not changed.
\Example
To transfer the system to the disk in drive~D, type \type{SYS D:}.
\end{DosEntry}

% TIME -- os/milton/command.c builtin_time, cmd_parse_time (hh:mm[:ss],
% 0..23 / 0..59; "Invalid time" then the unchanged time),
% cmd_format_time ("Current time is hh:mm:ss.cc").  Gates: test-datetime,
% test-command.
\begin{DosEntry}{TIME}{}{internal}
\Purpose Sets or displays the system time.
\Format \SyntaxLine{\kw{TIME} \opt{\ph{hh}\kw{:}\ph{mm}\opt{\kw{:}\ph{ss}}}}
\Type{Internal}
\Remarks
\Param{hh} is the hour, 0 through 23, on the 24-hour clock.

\Param{mm} is the minutes, 0 through 59.

\Param{ss} is the seconds, 0 through 59.  If you leave them out, the
seconds are set to 0.

With a valid time, \cmd{TIME} sets the system clock and displays the new
time.  With no parameter it displays the current time:
\texttt{Current time is} \ph{hh}\texttt{:}\ph{mm}\texttt{:}\ph{ss}%
\texttt{.}\ph{cc}, where \ph{cc} is hundredths of a second.

If the time is not valid, \cmd{TIME} displays \texttt{Invalid time}, leaves
the clock unchanged and displays the current time.
\Example
To set the clock to a quarter to two in the afternoon:
\Shot{time}
\end{DosEntry}

% TREE -- os/milton/xcopy_tree.c du_main(tree) (one path operand, default
% the root; /F; volume heading; "Path:", "Sub-directories:", ":" at column
% 16, "None", "Files:"; messages).  Shipped as APPS\TREE.COM.
% Gates: test-xcopy-tree, test-xcopy-tree-emu.
\begin{DosEntry}{TREE}{}{external}
\Purpose Displays the directories of a disk and, optionally, the files in
each.
\Format \SyntaxLine{\opt{\ph{d}\kw{:}}\opt{\ph{path}}\kw{TREE}
  \opt{\ph{d}\kw{:}}\opt{\ph{path}} \opt{\kw{/F}}}
\Type{External}
\Remarks
\Param{path} names the directory to start from.  If you leave it out,
\cmd{TREE} starts at the root directory.

\Param{/F} also lists the files in each directory.

\cmd{TREE} displays the heading \texttt{DIRECTORY PATH LISTING FOR VOLUME}
and the volume label.  Then, for the starting directory and each directory
below it, it displays \texttt{Path:} and the path, and
\texttt{Sub-directories:} followed by the names of the subdirectories, one
to a line, or \texttt{None}.  With \kw{/F}, \texttt{Files:} and the names
of the files follow, or \texttt{None}.

If the path is not a directory, \cmd{TREE} displays \texttt{Invalid path
(or file not found)}.  A parameter other than \kw{/F} displays
\texttt{Invalid parameter(s)}, and a second path displays \texttt{Invalid
number of parameters}.
\Example
To display the directories of the practice diskette:
\Shot{tree}
To list the files in the directory \file{SAMPLES}:
\Shot{treef}
\end{DosEntry}

% TYPE -- os/milton/command.c builtin_type (first token; open by name, no
% wildcard search; MSG-DOS-0003, MSG-DOS-0011; CR LF after the file).
% Gates: test-command, test-shell.
\begin{DosEntry}{TYPE}{}{internal}
\Purpose Displays the contents of a file.
\Format \SyntaxLine{\kw{TYPE} \opt{\ph{d}\kw{:}}\opt{\ph{path}}\ph{filename}\opt{\kw{.}\ph{ext}}}
\Type{Internal}
\Remarks
\Param{filename} names the file to display.  Wildcards are not allowed.

The file is displayed exactly as it is stored.  Use \cmd{TYPE} for text
files only; a program file contains codes that are not meant to be
displayed.

If the file does not exist, \cmd{TYPE} displays \texttt{File not found}.
With no file name it displays \texttt{Required parameter missing}.
\Example
To display the telephone list:
\Shot{type}
\end{DosEntry}

% VER -- os/milton/command.c builtin_ver (AH=30h: 3.30; "InitechDOS
% Version 3.30" after a blank line); spec/dos_banner.txt.  Gates:
% test-command, test-shell.
\begin{DosEntry}{VER}{Version}{internal}
\Purpose Displays the version number of InitechDOS.
\Format \SyntaxLine{\kw{VER}}
\Type{Internal}
\Remarks
\cmd{VER} displays a blank line and then the name and version of the
operating system.  Parameters are ignored.
\Example
\Shot{ver}
\end{DosEntry}

% VOL -- os/milton/command.c builtin_vol (drive A only, MSG-DOS-0004;
% extra parameters MSG-DOS-0010; MSG-DOS-0021/0022; read failure
% MSG-DOS-0009).  Gates: test-label-vol, test-dir-volume.
\begin{DosEntry}{VOL}{Volume}{internal}
\Purpose Displays the volume label of a disk.
\Format \SyntaxLine{\kw{VOL} \opt{\ph{d}\kw{:}}}
\Type{Internal}
\Remarks
\cmd{VOL} displays \texttt{Volume in drive A is} and the label, or
\texttt{Volume in drive A has no label}.  Create or change the label with
\cmd{LABEL}.

A drive other than A displays \texttt{Invalid drive specification}.
Anything after the drive displays \texttt{Too many parameters}.  If the disk
cannot be read, \cmd{VOL} displays \texttt{Access denied}.
\Example
\Shot{vol}
\end{DosEntry}

% XCOPY -- os/milton/xcopy_tree.c du_main (operands; /A /D:date /E /M /P /S
% /V /W; F/D question; names copied; "n File(s) copied"; exit codes 0, 1,
% 4, 5; cyclic and self-copy refusals; date/time kept via AH=57h; copies
% get archive set, source read-only kept).  Relative source from a
% subdirectory resolves against the root in the shipped build (no bead
% yet; WL-0119).  Shipped as APPS\XCOPY.COM.  Gates: test-xcopy-tree,
% test-xcopy-tree-emu, test-xcopy-tree-emu-cancel.
\begin{DosEntry}{XCOPY}{Extended Copy}{external}
\Purpose Copies groups of files, and the directories that contain them.
\Format \SyntaxLine{\opt{\ph{d}\kw{:}}\opt{\ph{path}}\kw{XCOPY}
  \opt{\ph{d}\kw{:}}\opt{\ph{path}}\opt{\ph{filename}\opt{\kw{.}\ph{ext}}}
  \opt{\ph{d}\kw{:}}\opt{\ph{path}}\opt{\ph{filename}\opt{\kw{.}\ph{ext}}}
  \opt{\kw{/A}\texttt{|}\kw{/M}} \opt{\kw{/D:}\ph{date}} \opt{\kw{/E}}
  \opt{\kw{/P}} \opt{\kw{/S}} \opt{\kw{/V}} \opt{\kw{/W}}}
\Type{External}
\Remarks
The first file specification names the source files.  It may be a file
name, a name with wildcards, or a directory, which means every file in it.

The second names the target.  If you leave it out, the files are copied to
the current directory.  If the target ends with a backslash, or is an
existing directory, the files are copied into it, and the directory is
created if necessary.  If the target does not exist and does not end with a
backslash, \cmd{XCOPY} asks:
\begin{TypeIn}
Does name specify a file name or directory name on the target
(F = file, D = directory)?
\end{TypeIn}
Press \type{F} to copy into a file of that name, or \type{D} to create a
directory.

\cmd{XCOPY} displays the name of each file as it copies it, and at the end
the number of files copied.  Each copy keeps the date and time of its
source and has its archive attribute set; a read-only source gives a
read-only copy.

\Param{/A} copies only files whose archive attribute is set, and leaves the
attribute unchanged.

\Param{/M} copies only files whose archive attribute is set, and clears the
attribute of each source file copied.  Use \kw{/M} to copy only the files
that have changed since the last copy.

\Param{/D:date} copies only files changed on or after \ph{date}, given as
for the \cmd{DATE} command.

\Param{/S} also copies the files in every subdirectory of the source,
creating the matching directories in the target.  Directories that hold no
files to copy are not created.

\Param{/E} used with \kw{/S} also creates the empty subdirectories.
\kw{/E} without \kw{/S} displays \texttt{Invalid parameter(s)}.

\Param{/P} asks before each file:
\ph{filename} \texttt{(Y/N?)}.  Press \type{Y} to copy the file or \type{N}
to skip it.

\Param{/V} reads back each copy and compares it with the source.

\Param{/W} displays \texttt{Press any key when ready to start copying
files} and waits, so that you can change diskettes first.

\cmd{XCOPY} will not copy a directory into one of its own subdirectories:
it displays \texttt{Cannot perform a cyclic copy}.  A file copied onto
itself displays \texttt{File cannot be copied onto itself}.

\cmd{XCOPY} returns exit code 0 when it copied files, 1 when no files were
found to copy, 4 for an error in the command line or a path, and 5 when a
file could not be read or written.

\draft{see bead initech-ktqk: run from a subdirectory, the shipped XCOPY finds a
relative source name in the root directory (XCOPY *.BAT in A:\bs SAMPLES
copies A:\bs AUTOEXEC.BAT); DOS 3.3 uses the current directory}
A source name without a path refers to the current directory.
\Example
The first command copies the text files into a new directory,
\file{BACKUP}; the second copies the directory \file{SAMPLES} to a new
directory, \file{SAVED}, the answer \type{D} being typed at the question:
\Shot{xcopy}
To be asked about each file, use \kw{/P}; here the answers are \type{Y},
\type{N} and \type{Y}:
\Shot{xcopyp}
\end{DosEntry}

